% !PRINTED 326
\weberpage{326}
Um die geometrische Bedeutung der singulären Punkte zu erkennen, denken wir uns die Function $f$ nach Potenzen der einen Variablen $x_3$ geordnet:

\weq{2}{f(x_1, x_2, x_3) = x_3^n f_0 + x_3^{n-1} f_1 + x_3^{n-2} f_2 + \cdots,}
worin $f_0, f_1, f_2, \dots$ binäre Formen der Variablen $x_1, x_2$ sind, deren Grad durch den Index angegeben ist, so dass $f_0$ eine Constante, $f_1$ eine lineare, $f_2$ eine quadratische Form ist u.\,s.\,f.

Wenn nun ein singulärer Punkt vorhanden ist, so können wir wegen der Invarianz dieser Eigenschaft annehmen, dass dieser Punkt in die Ecke $x_1 = 0$, $x_2 = 0$ des Coordinatendreiecks, die wir mit $\xi_3$ bezeichnen wollen, falle. Dann müssen $f_0$ und $f_1$ verschwinden, und wenn wir dann $x_1 = 0$ setzen, so erhält die Gleichung $f = 0$ den Factor $x_2^2$. Dies besagt, dass zwei der Schnittpunkte der Linie $x_1 = 0$ in dem singulären Punkte zusammenfallen. Die Linie $x_2 = 0$ kann aber jede durch den singulären Punkt hindurchgehende Gerade sein, und so folgt, dass der singuläre Punkt die Eigenschaft hat, dass \gesperrt{jede} durch ihn hindurchgehende Gerade die Curve dort in zwei zusammenfallenden Punkten schneidet. % [sic: the source prints $x_2=0$ here; the preceding derivation uses the cutting line $x_1=0$]

Aus diesem Grunde nennt man die singulären Punkte auch \gesperrt{Doppelpunkte}; womit aber nicht ausgeschlossen ist, dass auch mehr als zwei Schnittpunkte zusammenfallen können, wodurch die \gesperrt{höheren Singularitäten} entstehen.

An die Form der Gleichung (2) knüpfen wir noch einige in der Folge wichtige einfache Betrachtungen.

Wenn der Punkt $\xi_3$ auf der Curve $f$ liegt, aber kein singulärer Punkt ist, so muss die Constante $f_0$ verschwinden und $f_1 = 0$ ist die Gleichung der \gesperrt{Tangente} der Curve in diesem Punkte. Nehmen wir diese Tangente für die Linie $x_1 = 0$, so lautet die Gleichung der Curve, wenn wir einen constanten Factor gleich 1 annehmen:

\weq{3}{f = x_3^{n-1} x_1 + x_3^{n-2} f_2 + \cdots = 0.}
Wenn die quadratische Form $f_2$ durch $x_1$ theilbar ist, so erhält die Gleichung (3) für $x_1 = 0$ den Factor $x_2^3$ und der Punkt $\xi_3$ zählt als dreifacher Schnittpunkt von $x_1 = 0$ mit der Curve. Ein solcher Punkt heisst ein \gesperrt{Wendepunkt} oder \gesperrt{Inflexionspunkt} der Curve $f$, und $x_1 = 0$ heisst die \gesperrt{Wendetangente}. Die Wendepunkte gehören nicht zu den singulären Punkten.
